\documentclass[12pt,a4paper]{article}
\usepackage[utf8]{inputenc}
\usepackage[T1]{fontenc}
\usepackage{amsmath}
\usepackage{amssymb}
\usepackage{amsfonts}
\usepackage{mathtools}
\usepackage{physics}
\usepackage{siunitx}
\usepackage{graphicx}
\usepackage{xcolor}
\usepackage{hyperref}
\usepackage{geometry}
\usepackage{microtype}
\usepackage{lmodern}
\geometry{margin=2.5cm}
\hypersetup{colorlinks=true,linkcolor=blue,urlcolor=blue}

\title{Practical-1.pdf}
\author{Jack Ma}
\date{07 mai 2026}

\begin{document}
\maketitle

\section*{Practical-1.pdf}
\begin{figure}
\includegraphics[alt={},max width=\textwidth]{https://cdn.mathpix.com/cropped/a7d4e9ef-534e-4c7f-8498-c99e82184750-1.jpg?height=405&width=997&top_left_y=324&top_left_x=529}
\captionsetup{labelformat=empty}
\caption{image analysis}
\end{figure}

MPhil in Data Intensive Science

Thursday 7rd May 2026

Minor: Image Analysis
- Problem Sheet 1: Image transforms and classical image processing

Dr A. Biguri

In this practical session you will learn to perform basic tasks with images, and to use some of the standard image transforms.

Attempt all questions.
We will work through the questions during the practical session on the date indicated above.

Solutions to these questions will be available on Moodle after such practical session.

\section*{Module 1: Image and noise manipulation}

Question 1: Using ( skimage ), load and display a couple of images, e.g. coins and cat. Inspect the size, type, number of channels of each of them. e

Question 2: Lets try to segment the coins. As they are brighter than the background, thresholding may work.

Question 3: The thresholding is close, but not sufficient. Use morphological operation to get circles for each coin. https://scikit-image.org/docs/stable/api/skimage.morphology.html

Question 4: Lets try to get rid of things that are not coins. For that, we need to be able to select individual things. How?

Question 5: Now that we have individual items, lets try to get rid of the ones that are not a coin. We can do that by computing each region properties.
[https://scikit-image.org/docs/stable/api/skimage.measure.html\# skimage.measure.regionprops

Question 6: Lets pick one of them. E.g. the biggest one.

Question 7: Use the label to extract the coin only from the original image. Lets also remove unnecessary blank space.

Question 8: Lets give the cat something to look at. Lets put the coin in front of its face, rotated 30 clockwise. In front of its eyes is approximately at [200,350]. Do not use inbuilt transform functions, but code it yourself.

Question 9: Can you do the same thing using inbuilt functions?


\end{document}
